\section{Lecture 5, Part 1: Improving Retrieval Quality}

\begin{frame}
\vfill
\centering
{\color{RLBlue}\large Lecture 5}

\vspace{0.3em}
{\color{RLBlue}\Large\bfseries Part 1}

\vspace{0.6em}
{\Large Improving Retrieval Quality}
\vfill
\end{frame}

\begin{frame}{Learning Goals}
\begin{itemize}
  \item Say why a second ranking stage exists, and what a cross-encoder can do that a first-stage retriever cannot.
  \pause
  \item Combine two ranked lists with reciprocal rank fusion, by hand.
  \pause
  \item Choose between HyDE, multi-hop retrieval and retrieval-on-demand by naming the failure each one repairs (Gao et al., 2022; Asai et al., 2023).
\end{itemize}
\end{frame}

% ─── Reranking ──────────────────────────────────────────────────────────────
\begin{frame}{Why Isn't First-Stage Retrieval Enough?}
\begin{itemize}
  \item \textbf{Why:} the first stage must score \emph{every} chunk in the corpus, so it is built for speed. That budget buys a scoring function that looks at the query and the chunk \emph{separately} and compares the two summaries.
  \pause
  \item Scoring them separately is what makes the index possible -- chunk scores can be precomputed -- and it is also what loses the interaction. ``Reduced working time'' and ``employees under eighteen'' both match Article~183 strongly on their own; only reading them \emph{against} the article reveals it states no limit.
  \pause
  \item \textbf{So use two stages.} Retrieve $50$ candidates cheaply, then rescore those $50$ expensively. You pay the expensive model $50$ times per query instead of $7{,}368$ times.
  \pause
  \item \textbf{The asymmetry is the whole design:} recall is the first stage's job and cannot be recovered later; precision is the second stage's job and is cheap to buy.
\end{itemize}
\end{frame}

\begin{frame}{The Retrieve-then-Rerank Funnel}
\centering
\begin{tikzpicture}[
  font=\scriptsize,
  band/.style={draw=RLBlue, fill=RLGray, rounded corners=2pt, inner sep=4pt,
               minimum height=7mm, align=center},
  lbl/.style={font=\scriptsize, inner sep=2pt, align=center}
]
  \node[band, minimum width=2.9cm] (all)  at (0,0)    {$7{,}368$ chunks};
  \node[band, minimum width=2.2cm] (top50) at (4.3,0)  {top $50$};
  \node[band, minimum width=1.6cm] (top5)  at (8.3,0)  {top $5$};
  \node[band, minimum width=1.9cm, draw=RLGreen] (gen) at (11.2,0) {generator};
  \draw[->] (all)   -- node[lbl,above=3.2mm] {first stage} (top50);
  \draw[->] (top50) -- node[lbl,above=3.2mm] {cross-encoder} (top5);
  \draw[->] (top5)  -- (gen);
\end{tikzpicture}

\vspace{0.8em}
\small
\begin{itemize}
  \item \textbf{The cost asymmetry is the point.} The first stage scores all $7{,}368$ chunks, so it must be cheap. The cross-encoder scores only $50$, so it can afford a forward pass each.
  \pause
  \item \textbf{Recall is set at the first arrow:} if the gold article is not among the $50$, nothing downstream can find it.
  \pause
  \item \textbf{Precision is set at the second:} the reranker's job is to get the gold article into the $5$ the generator actually reads.
\end{itemize}
\end{frame}

\begin{frame}{Two Ways to Score a Query Against a Chunk}
\begin{rldefinition}
\textbf{First-stage retriever} -- encodes the query and each chunk \emph{independently} and compares the two summaries. Chunk side is precomputed once, so scoring the whole corpus is cheap. Lecture~6 builds one.

\smallskip
\textbf{Cross-encoder} -- feeds query and chunk to the model \emph{together}, as one sequence, and outputs a single relevance score. Nothing can be precomputed: every (query, chunk) pair costs a forward pass.

\smallskip
\emph{Here:} the first stage proposes Art.~183 and Art.~415; the cross-encoder reads ``how many hours may a 17-year-old work?'' \emph{alongside} each one and notices only Art.~415 contains a number.
\end{rldefinition}
\pause
\begin{itemize}
  \item Nogueira \& Cho (2019) introduced this reranking design for passage ranking, and it remains the standard second stage.
\end{itemize}
\end{frame}

% ─── RRF ────────────────────────────────────────────────────────────────────
\begin{frame}{Reciprocal Rank Fusion}
\begin{rldefinition}
\textbf{Reciprocal rank fusion (RRF)} -- score each document by its reciprocal rank summed over every list it appears in, with $k = 60$ by convention (Cormack, Clarke \& Buettcher, 2009):
\vspace{-0.4em}
\[
  \mathrm{RRF}(d) \;=\; \sum_{\ell \in \text{lists}} \frac{1}{k + \mathrm{rank}_{\ell}(d)} .
\]
\vspace{-1.2em}

\emph{Here:} fuse a keyword ranking with Lecture~6's dense ranking.
\end{rldefinition}
\pause
\begin{itemize}
  \item \textbf{Why ranks, not scores:} two retrievers return numbers on incomparable scales -- a similarity in $[0,1]$ against an unbounded term-weight sum. Ranks need no calibration.
  \pause
  \item \textbf{Why the $+\,k$:} it flattens the top, so first rather than third is worth little. Agreement across lists beats dominance within one.
\end{itemize}
\end{frame}

\begin{frame}{Worked Example: Fusing Two Rankings}
\small
Two retrievers rank four candidate articles for \emph{``how many hours may a 17-year-old work?''}
\smallskip

\centering
\footnotesize
\begin{tabular}{@{}llcccl@{}}
& & \textbf{List A} & \textbf{List B} & & \\
\textbf{Article} & & rank & rank & \textbf{RRF score} & \\
\hline
Art.~258 & remuneration for reduced hours & $1$ & --      & $1/61 = 0.0164$ & \\
Art.~183 & reduced hours, full pay        & $2$ & $3$     & $1/62 + 1/63 = 0.0320$ & \\
\textcolor{RLGreen}{Art.~415} & \textcolor{RLGreen}{16--18: 36 h/wk} & $3$ & $1$ & \textcolor{RLGreen}{$1/63 + 1/61 = 0.0323$} & \\
Art.~182 & normal working time            & $4$ & $2$     & $1/64 + 1/62 = 0.0318$ & \\
\end{tabular}
\pause

\vspace{0.5em}
\begin{itemize}
  \item \textbf{Art.~258 was ranked first by List A and finishes last.} It appears in no other list, so it scores once. One first place is worth about \emph{half} of two middling places.
  \pause
  \item \textbf{Art.~415 was never first in List A and wins.} Two retrievers that disagree about the order still agree it is relevant, and that agreement is what RRF rewards.
\end{itemize}
\end{frame}

% ─── HyDE ───────────────────────────────────────────────────────────────────
\begin{frame}{Why Do We Need HyDE?}
\begin{itemize}
  \item \textbf{Why:} queries and documents are written in different registers. A student asks \emph{``how many hours may a 17-year-old work?''}; the statute says \emph{``reduced working time is established for employees aged sixteen to eighteen''}. They share almost no words.
  \pause
  \item \textbf{HyDE} (Gao et al., 2022) closes the gap from the query side: ask a model to \emph{write} a hypothetical answer, then retrieve using that text instead of the question. The invented answer is in the register of the corpus even when its facts are wrong.
  \pause
  \item \textbf{The counter-intuitive part:} the hypothetical document does not need to be correct. It is never shown to the user and never cited -- it is only a better-shaped query. Gao et al.\ show this works with no relevance labels at all.
  \pause
  \item \textbf{When it hurts:} if the model's invention is confidently off-topic it drags retrieval with it. Fusing the HyDE ranking with the original-query ranking -- by RRF -- keeps the upside and caps the downside.
\end{itemize}
\end{frame}

% ─── Multi-hop ──────────────────────────────────────────────────────────────
\begin{frame}{Multi-Hop Retrieval: Following a Cross-Reference}
\begin{itemize}
  \item Some questions cannot be answered by any single chunk, however well you cut the corpus. The under-18 question is one: Article~183 establishes \emph{that} reduced hours exist; Article~415 says \emph{how many}.
  \pause
  \item \textbf{Multi-hop retrieval} runs retrieval again using what the first pass returned -- the agent loop of Lecture~1, with \texttt{search} as the tool. Hop~1 finds Art.~183; the thought ``this names a category but no number'' drives hop~2.
  \pause
  \item \textbf{This is why an agent beats a pipeline here.} A fixed pipeline retrieves once. Nothing in one pass of Art.~183 tells a pipeline that a second search is needed.
  \pause
  \item \textbf{Free ground truth:} the corpus stores \texttt{cross\_references}, so article~A citing article~B gives a two-hop query whose gold answer set is known without an annotator. Build your multi-hop evaluation from that field.
\end{itemize}
\end{frame}

\begin{frame}{Self-RAG: Deciding Whether to Retrieve At All}
\begin{itemize}
  \item Everything so far retrieves once per query, always. Both halves of that are wrong: some queries need several passes, and some need none -- retrieving for ``summarise what you just told me'' only adds noise.
  \pause
  \item \textbf{Self-RAG} (Asai et al., 2023) trains the model to emit \textbf{reflection tokens} alongside its text, drawn from an expanded vocabulary. A \texttt{[Retrieve]} token decides on demand whether to search at all.
  \pause
  \item Three further tokens critique what came back and what was written with it: whether the passage is \emph{relevant}, whether the output is \emph{supported} by it, and whether the answer is \emph{useful}.
  \pause
  \item \textbf{Why this belongs in a legal agent:} ``supported by the retrieved passage'' is exactly the property Lecture~8 grades. Self-RAG makes the model state that judgement in the trace, where a harness can read it -- rather than leaving it to be inferred from the final sentence.
\end{itemize}
\end{frame}

\begin{frame}{Fusion-in-Decoder: Using More Passages Than Fit}
\begin{itemize}
  \item Pasting $k$ chunks into one prompt costs attention quadratic in the total length, which is what caps $k$ at a handful.
  \pause
  \item \textbf{Fusion-in-Decoder} (Izacard \& Grave, 2021) encodes each retrieved passage \emph{separately} -- each paired with the question -- then concatenates the encoder outputs and lets the decoder attend across all of them at once.
  \pause
  \item Encoding becomes linear in the number of passages and fully parallel, so the same model can read far more: they report Fusion-in-Decoder at $100$ passages reaching $51.4\%$ exact match on Natural Questions, against $47.5\%$ for the RAG-Sequence model of Lewis et al.
  \pause
  \item \textbf{The caveat for this module:} it needs an encoder--decoder you can modify, so it is unavailable behind a hosted chat API. Know it as the reason ``just paste more passages'' is not the only answer to low recall.
\end{itemize}
\end{frame}

\begin{frame}{Why This Matters: recall@\texorpdfstring{$k$}{k} Is the Score}
\begin{itemize}
  \item Sprint~2 reports recall@$k$ against the Sprint~1 baseline. Each technique in this lecture moves a different term, and confusing them wastes a sprint:
  \pause
  \item \textbf{Chunking} raises the ceiling -- it changes what \emph{can} be retrieved. \textbf{Fusion and HyDE} raise first-stage recall -- more gold articles inside the $50$. \textbf{Reranking} raises precision -- it cannot add a gold article, only promote one already there.
  \pause
  \item \textbf{Multi-hop} changes what counts as answerable at all, and needs its own evaluation set built from \texttt{cross\_references}.
  \pause
  \item \textbf{Report them as an ablation, not a story.} One row per configuration, recall@$5$ and recall@$50$ side by side. The gap between those two columns tells you whether to work on the retriever or on the reranker.
\end{itemize}
\end{frame}

\begin{frame}{Summary}
\begin{itemize}
  \item \textbf{RAG} replaces remembering with reading: retrieve, augment, generate -- and every claim can then name its source.
  \pause
  \item \textbf{Chunking sets the ceiling.} A clause split in half is unretrievable by any ranking function, and legal text hands you the parent--child structure for free.
  \pause
  \item \textbf{Two stages, two jobs:} the first stage owns recall and cannot be corrected later; the cross-encoder owns precision and is cheap because it only ever sees $50$ candidates.
  \pause
  \item \textbf{RRF} fuses rankings without calibrating scores, and rewards agreement across retrievers over dominance within one.
  \pause
  \item \textbf{HyDE} repairs vocabulary mismatch, \textbf{multi-hop} repairs answers split across articles, and \textbf{Self-RAG} repairs retrieving when you should not have.
\end{itemize}
\end{frame}

\begin{frame}{Reading}
\begin{itemize}
  \item Lewis, Perez, Piktus, Petroni, Karpukhin, Goyal, K\"uttler, Lewis, Yih, Rockt\"aschel, Riedel \& Kiela (2020). \textit{Retrieval-Augmented Generation for Knowledge-Intensive NLP Tasks.} NeurIPS. --- parametric and non-parametric memory.
  \pause
  \item Izacard \& Grave (2021). \textit{Leveraging Passage Retrieval with Generative Models for Open Domain Question Answering.} EACL. --- Fusion-in-Decoder.
  \pause
  \item Gao, Ma, Lin \& Callan (2022). \textit{Precise Zero-Shot Dense Retrieval without Relevance Labels.} --- HyDE.
  \pause
  \item Asai, Wu, Wang, Sil \& Hajishirzi (2023). \textit{Self-RAG: Learning to Retrieve, Generate, and Critique through Self-Reflection.} ICLR 2024.
  \pause
  \item \textbf{Next lecture:} what the first stage actually is --- bi-encoders, contrastive training, and why BM25 is a baseline worth beating.
\end{itemize}
\end{frame}
